% ID: SEL-D0050-OB-E01
% ESTADO: cotejada. Segunda pasada visual y borrador compilado revisados.
% Fuente: Selectividad/sel_2013_matematicas_aplicadas/junio/titular/EXAMEN 6 -ANDALUCÍA.pdf
% SHA-256: 1f0983e351b8435c5b6c7b68b05fe2a64034faf77a85a1f9480407aa0f2a88aa
% Página: 2; opción B; ejercicio original 1.

Un fabricante de tapices dispone de 500 kg de hilo de seda, 400 kg de hilo de plata y 225 kg de hilo de oro. Desea fabricar dos tipos de tapices: A y B. Para los del tipo A se necesita 1 kg de hilo de seda y 2 kg de hilo de plata, y para los del tipo B, 2 kg de hilo de seda, 1 kg de hilo de plata y 1 kg de hilo de oro. Cada tapiz del tipo A se vende a 2000 euros y cada tapiz del tipo B a 3000 euros. Si se vende todo lo que se fabrica,
\begin{enumerate}[label=\alph*)]
\item \textbf{(2 puntos)} \textquestiondown{}cuántos tapices de cada tipo ha de fabricar para que el beneficio sea máximo y cuál es ese beneficio?
\item \textbf{(0.5 puntos)} \textquestiondown{}Qué cantidad de hilo de cada clase quedará cuando se fabrique el número de tapices que proporciona el máximo beneficio?
\end{enumerate}
